\section{Autonomy: Always-On Operation, Night Shift, Self-Improvement, Desktop Control and Calls}
\label{sec:autonomy}

The features in this section move \hivi{} from a tool the user operates to an agent that works while the user is away. Each is a client of the security base (\sect{sec:security}); together they are the reason that base is built early.

\subsection{Always-on \hivi{} and phone control}

The goal: from the user's phone, anywhere, ``\hivi{}, until I reach home make me $X$, $Y$ and $Z$''---the work happens in the background and is done on arrival; routines run on their own; the assistant messages the user when something needs them.

\paragraph{What is physically possible} A shut-down laptop runs nothing; everything else is possible (Table~\ref{tab:laptop}). The practical set-up is a laptop kept plugged in and set to sleep rather than shut down, with wake-on-network enabled; \hivi{} keeps it awake only while jobs run.

\begin{table}[t]
\caption{What \hivi{} can do in each laptop state}
\label{tab:laptop}
\centering\footnotesize
\begin{tabularx}{\columnwidth}{@{}L{30mm}Y@{}}
\toprule
State & Capability \\
\midrule
On (lid closed is fine, plugged in) & Everything: tasks, routines, messages \\
Asleep & The phone can wake it (Wake-on-LAN; Ethernet best, Wi-Fi depends on the laptop), then it works \\
Shut down & Nothing runs; tasks queue and start the moment it boots \\
Optional later: small always-on box & Holds the queue and wakes the laptop; can run tiny models for quick answers; heavy work waits for the laptop \\
\bottomrule
\end{tabularx}
\end{table}

\paragraph{Phone control} A private connection through Tailscale (a WireGuard~\cite{wireguard} overlay between phone and laptop) so \hivi{} is reachable from anywhere without opening a port to the internet; the web UI as an installable phone app (PWA) with the same chats, memory and models, by text or voice; optionally a Telegram bot linked to the user's account alone; push notifications (``Done: 3 files ready'', ``Needs your OK'', ``Failed: here's why''); and login with device approval (\phase{SEC-3}), because anyone reaching \hivi{} could run code on the laptop.

\paragraph{Background jobs} A SQLite job queue: every task from the phone becomes a job with status (queued, running, needs approval, done, failed), progress and results. Jobs survive a closed browser, lost signal, a crash (resuming from the last checkpoint) and a reboot (\hivi{} starts with the OS as a background service). The runner keeps the laptop awake only while a job runs, uses the same engines as chat (search, agents, image and voice generation), and lands results on the job's page on the phone.

\paragraph{``Almost alive'': routines and watchers} Scheduled routines (``every morning at 8, summarise new data-science job postings''; ``every Sunday, back up my memories''); watchers that check a page, folder or repository on a schedule and act or notify on change; proactive messages only when something matters---a job finished, a watcher fired, an approval is needed---never spam.

\paragraph{Safety} Remote control is power over the laptop, so the phone-side view of the shared gate applies: reading, searching, writing inside the workspace, generating images and voice, and running sandboxed code proceed without asking; deleting or overwriting files outside the workspace, sending or posting anything, and spending, installing or changing system settings need approval on the phone. A kill switch on the phone stops all jobs instantly; the activity log shows everything done while the user was away; per-job budgets (time, model calls) ensure nothing runs forever.

\paragraph{Phases} \phase{LIVE-1}: job queue and background runner, starts with the OS, keep-awake, job status page. \phase{LIVE-2}: phone access (Tailscale, PWA, push). \phase{LIVE-3}: approvals from the phone, kill switch, activity-log viewer. \phase{LIVE-4}: routines and watchers. \phase{LIVE-5}: Telegram bot, voice from the phone. \phase{LIVE-6}: wake from sleep, optional always-on box.

\subsection{Night shift}

While the user sleeps, \hivi{} does the maintenance and heavy work that would slow it down during the day, then leaves a short morning report with anything needing a decision. It runs in a window the user sets (for example 2--6~a.m.), only when plugged in and with normal GPU temperature; the user always wins---starting to use the PC or messaging from the phone pauses night work instantly. The night's job list, in priority order: backup snapshot; health check (disk space, log clean-up, engine self-test); memory tidy (merge duplicates, flag outdated facts for review---never deleting on its own); re-index changed files and rebuild embeddings after an embedding-model change; regression exams on the current models; the model scout (candidates downloaded and examined); every skill's tests, flagging broken ones; research prep (pre-searching topics from tomorrow's open tasks and watched subjects so morning answers are instant); and long jobs the user queued (overnight experiments, agent runs, image batches). Limits: a time, energy and disk budget per night, with checkpoints so a job continues the next night; anything that sends, spends or changes files outside the workspace is not done at night and waits in the morning report. The report is one card in the app and a phone notification: what ran, what failed, backup status, and the decisions waiting (``switch to this model?'', ``merge these 3 memories?'', ``keep this new skill?''), each with one-tap buttons.

\paragraph{Phases} \phase{NIGHT-1}: window, plug-in and temperature checks, pause when active; backups, health checks, re-index; morning report. \phase{NIGHT-2}: memory tidy suggestions, exams and scout, skill tests, research prep, queued long jobs.

\subsection{Self-improvement: models, skills, workflows and blueprints}

Once a day \hivi{} checks whether something better exists---a new model, a better way to do a skill, a better workflow---tests it on the user's own machine, and asks: ``I found $X$; it scored better than what I use now on these tests. Want me to switch?''

\paragraph{The core rule} Nothing changes unless it wins on \hivi{}'s own exams (P7). Public leaderboards do not know the user's hardware, tasks or settings, so every role and department has a \emph{local exam}: a fixed set of tasks with automatic scoring, and a candidate replaces the incumbent only if it scores better here at a speed the user accepts. Table~\ref{tab:exams} lists the exams; \sect{sec:evaluation} returns to them as the paper's evaluation methodology.

\begin{table}[t]
\caption{Local exams}
\label{tab:exams}
\centering\footnotesize
\begin{tabularx}{\columnwidth}{@{}L{19mm}YL{24mm}@{}}
\toprule
Exam & What it tests & Scored by \\
\midrule
Coding & Fix / implement tasks with hidden tests (the agent-company benchmark suite) & Tests pass \\
Planning / reviewing & Plans and reviews checked against known answers & Rubric + code checks \\
Chat and personality & Fixed conversations; tone and honesty rules (\sect{sec:soul}) & Rubric + the user's past thumbs \\
Search answers & Questions with known answers and sources & Correct facts, correct citations \\
Uncensored builds & The user's refusal test set; drift vs.\ the base model & Refusal count; quality kept \\
Images / voice & Fixed prompts; speed and VRAM & The user's side-by-side picks + speed \\
Speed & Tokens/s, time to first token, RAM and VRAM on this PC & Measured \\
\bottomrule
\end{tabularx}
\end{table}

\paragraph{Daily model scout} (1) Scan, as a scheduled routine: new and trending models on Hugging Face, release notes of model families in use, engine updates. (2) Filter to what fits: 32~GB RAM and 8~GB VRAM; MoE with at most about 8B active parameters; no quantisation worse than Q5 for MoE; an acceptable licence; known uploaders; GGUF or safetensors only (no pickle files); uncensored builds only with a measured refusal and divergence report. (3) Shortlist one to three candidates per role from published results. (4) Download overnight within a disk budget and run the exams on idle time. (5) Report: ``Candidate vs.\ current: +12\% coding tests passed, same speed, 3~GB bigger'', with Switch / Try for a week / Skip buttons in the app or on the phone. (6) Switch safely: the old model stays installed; automatic rollback if the new one fails exams again or the user says ``go back''; unused old models are offered for deletion, never deleted silently.

\paragraph{Skills, workflows and blueprints} Skills improve the same way: weekly re-checks of the sources behind each skill; a new version is built beside the old one and must pass the skill's tests and beat it on quality or speed; failures in real use trigger a revision first. Workflows---how \hivi{} does multi-step work: agent pipelines, the search loop, HAIE's sampling rules, prompt templates, retry budgets, the number of worker slots---improve the same way: \hivi{} proposes variants (a different prompt for the test author, four worker slots instead of three, one more search round), runs them on the exams in the background, proposes a winner with evidence, and records why in the decision log. A guardrail: variants may change prompts, settings and model assignments, never the safety rules, approval rules or exam answers.

\paragraph{The user approves every upgrade} Default: ask before every switch. Optional: auto-apply small wins (settings, prompts) and ask only for model and skill swaps. A weekly ``what got better'' summary with exam scores before and after and what was rolled back. One switch pauses all self-improvement.

\paragraph{\hivi{} improves \hivi{}} The assistant points its own agent company at its own repository: it fixes bugs from its error log, speeds up slow parts and adds small features the user asks for---always on a separate copy, always gated by tests, always with approval. Ideas come from repeated errors in the error log and slowdowns in the speed history (\sect{sec:operations}), failed exams, patterns in thumbs-down feedback, and direct requests. The loop: an issue with evidence (log lines, timings) and a proposed fix; work on a separate branch in a separate folder (a git worktree), never on the running \hivi{}; gates---the full test suite, the exams, and a new test that proves the fix; the user's review of summary, diff and test results; apply safely---snapshot, apply, restart, health check; automatic rollback if the health check or exams fail, or on request within 24 hours. A \emph{protected zone} the assistant may propose changes to but never make: the security base (login, approvals, vault, sandbox, activity log), so it cannot loosen its own limits; exam answers and existing tests, which it may add to but never delete or weaken (the gate checks); and the self-improvement rules and the rollback and backup code themselves, so a bad change can always be undone. Limits: one change at a time, small changes only (a set number of files and lines), a few per week; anything bigger becomes a plan for the user to approve first.

\paragraph{Phases} \phase{UP-1}: local exam suites per role and a scoreboard. \phase{UP-2}: daily scout (scan, filter, shortlist, report). \phase{UP-3}: overnight download and exam runs, Switch / Try / Skip, rollback. \phase{UP-4}: skill self-improvement. \phase{UP-5}: workflow and blueprint variants on exams. \phase{UP-6}: engine update checks with regression tests. \phase{SELF-1}: self-coding, read-only (issues with evidence and proposed patches the user applies). \phase{SELF-2}: applies approved fixes under branch, gates, protected zone, backup, rollback. \phase{SELF-3}: small features from requests; optional auto-apply for tiny fixes that pass every gate, with notification.

\subsection{Desktop control}

\hivi{} uses desktop applications the way the user does---a spreadsheet, the file manager, settings, any app---to finish tasks that have no connector or API. The most reliable method is used first: a direct way (connector, shell, application automation, a file library) whenever one exists; the accessibility tree (buttons, menus and fields by name, e.g.\ via pywinauto) for most apps, which is text rather than pixels and works well with small models; and screenshots with the vision model for apps that do not expose their controls, which works almost anywhere but is slower and less reliable.

This is the strictest feature in \hivi{}. Off by default and enabled per task. An app allow-list, never including password managers, banking, security settings or \hivi{}'s own security settings. Watch mode first: every step shown live and needing approval; later only risky steps ask. Stop instantly with a key on the PC and the kill switch on the phone. Screen content is untrusted (P6). It never fights the user for the mouse: it works only when invited or when the user is away, and later in a separate virtual machine so it never touches the user's session. Screenshots stay in the run log only and are deleted after a few days. A procedure that worked is saved as a skill, so next time is faster and more reliable; once \phase{DESK-3} is in place, phone jobs can include ``while I'm out, do $X$ in that app''.

\paragraph{Phases} \phase{DESK-1}: read-only---look at an app (tree plus screenshot) and guide the user step by step. \phase{DESK-2}: actions in allowed apps with approval per step, live view, stop key. \phase{DESK-3}: longer unattended tasks from the phone in a separate virtual machine; procedures saved as skills.

\subsection{Phone calls}

The final stage: \hivi{} calls the user when something urgent comes up, and the user can call \hivi{} and just talk. It builds on voice mode (\sect{sec:capabilities}) and phone control. Three ways to make the call (Table~\ref{tab:calls}); the app call is built first because it is fully private and free and works anywhere with data.

\begin{table}[t]
\caption{Three ways to place a call}
\label{tab:calls}
\centering\footnotesize
\begin{tabularx}{\columnwidth}{@{}L{17mm}YL{22mm}@{}}
\toprule
Way & How it rings & Privacy / cost \\
\midrule
App call (first) & The phone app rings through a high-priority push; audio phone $\leftrightarrow$ laptop over Tailscale (WebRTC) & Fully private, free \\
Private phone line & A SIP app on the phone registered to a small phone server on the laptop (Asterisk / FreeSWITCH) over Tailscale & Private, free; rings like a real call even with the app closed \\
Real phone number & A telephony provider gives \hivi{} a number; call audio streams to the laptop & Not local: the provider carries the audio; number fee plus per-minute cost \\
\bottomrule
\end{tabularx}
\end{table}

\paragraph{\hivi{} calls the user, urgent only} What counts as urgent is set by the user as rules: a job failed that blocks something due today, an approval needed before a deadline, a watcher marked ``call me''. An escalation ladder: notification first, then a call if there is no answer in $N$ minutes and it is urgent. Quiet hours and a do-not-disturb switch with chosen exceptions. A rate limit of a few calls a day, never repeating the same alert. On the call, \hivi{} says why in one sentence; the user can talk it through, approve or deny by voice (``yes, send it'' / ``no, wait till I'm home''), or say ``text me the details''. A missed call leaves a voicemail-style summary.

\paragraph{The user calls \hivi{}} The same voice conversation mode: interrupt, ask anything, give tasks that go into the job queue. The assistant knows the user's memory, running jobs and routines, so ``how's that job going?'' works. The transcript and anything done are saved to the activity log.

\paragraph{Safety and GPU budget} Caller ID can be faked, so on a real phone number \hivi{} asks for a spoken PIN or passphrase before doing anything beyond chatting; app and private-line calls are already tied to the user's devices. Voice approvals follow the phone rules; risky actions are read back and need an explicit ``yes''. The kill switch works by voice. A live call needs speech-to-text, the chat model and text-to-speech at once; a heavy job using the GPU is paused for the call and resumed after.

\paragraph{Phases} \phase{CALL-1}: app call, using voice mode. \phase{CALL-2}: \hivi{} calls the user (urgent rules, ladder, quiet hours, rate limits, missed-call summary). \phase{CALL-3}: voice approvals and voice kill switch. \phase{CALL-4}: private phone line. \phase{CALL-5}: optional real phone number with spoken PIN.
